\documentclass[10pt,a4paper]{amsart}
\usepackage[margin=19mm]{geometry}
\usepackage{amsmath,amssymb,mathtools}
\usepackage[scr=rsfs,cal=euler]{mathalfa}
\usepackage{xfrac}
\usepackage[unicode,hidelinks,bookmarks=false]{hyperref}
\newcommand{\ie}{i.e.\ }
\newcommand{\cf}{cf.\ }
\newcommand{\resp}{respectively\ }
\newcommand{\Subsection}{\subsection}
\providecommand{\nobreakdash}{\nobreak\hbox{-}}
\setlength{\emergencystretch}{2em}
\multlinegap0pt
\usepackage{amssymb}
\usepackage{mathtools}
\newtheorem{lemma}{Lemma}
\newtheorem{theorem}{Theorem}
\newcommand{\rank}{\operatorname{rank}}
\newcommand{\tr}{\operatorname{tr}}
\title{No involutions in the missing Moore graph}
\author{Yawara Ishida}
\date{Source: arXiv:2606.29183v2 (CC BY 4.0)}
\begin{document}
\maketitle

\noindent\emph{Source and licence.} No involutions in the missing Moore graph. Version arXiv:2606.29183v2 (CC BY 4.0), distributed by arXiv under the Creative Commons Attribution 4.0 International licence, http://creativecommons.org/licenses/by/4.0/, as recorded in the arXiv licence metadata for this exact version and shown on the abstract page https://arxiv.org/abs/2606.29183v2. Full licence notice: \texttt{licenses/arxiv-2606.29183-CC-BY-4.0.txt}.\par\smallskip

\noindent\textit{Editorial scope.} Counted unit: the article's theorem thm:local (source0.tex lines 371--380). The article's complete original proof of it is delivered with it (source0.tex lines 382--418, one \texttt{proof} environment of 37 lines). No mathematics is written, repaired or re-derived: the statement and the proof are the article's own lines, in the article's own notation, and no retained line is substituted or re-flowed.  Retained uncounted as the source-local context the proof needs: lines 271--271; lines 313--319; lines 367--369.  Nothing was omitted from any retained span, so the package carries no omission record.  No retained line contains a citation, so the wrapper carries no bibliography and no bibliography entry.  No retained line contains a figure environment, an image-inclusion command or drawing code, so no image is delivered and the package declares no asset dependency.  Editorial formatting only, in addition: an AMS \texttt{amsart} preamble that restates the article's own declarations and macros used by the retained lines, namely the environments \texttt{lemma}, \texttt{theorem} on separate counters, together with the standard packages those lines need.  The retained environments print in this document as Lemma 1, Theorem 1 in order of appearance, whereas the article prints the same items with the numbering of its own full document.  The complete native source of the article as distributed in the arXiv e-print for this exact version is bundled unchanged as \texttt{sources/203910/source0.tex}.
\section{\texorpdfstring{A trace formula for $p$-permutation lattices}{A trace formula for p-permutation lattices}}\label{sec:padic}

\begin{lemma}\label{lem:trace-brauer}
Let $N$ be a $p$-permutation $OC_p$-lattice.  Then
\[
 \tr\left(x\mid K\otimes_{O}N\right)=\dim_{\kappa}N(C_p),
\]
where the integer on the right is regarded as an element of $K$.
\end{lemma}

\begin{equation}\label{eq:MC-basis}
 M(C_p)\cong \kappa^{\mathcal F}.
\end{equation}

\begin{theorem}\label{thm:local}
In the situation above, let $E\in\operatorname{End}_O(M)$ be an idempotent commuting with the action of $C_p$,
set $L=EM$,
and let $E[\mathcal F]$ denote the principal submatrix of the matrix of $E$ in the standard basis of $M$ on the rows and columns indexed by $\mathcal F$.
Then
\[
 \tr\bigl(x\mid K\otimes_O L\bigr)=\rank_{\kappa}\bigl(E[\mathcal F]\bmod \mathfrak m\bigr),
\]
the integer on the right being regarded as an element of $K$.
\end{theorem}

\begin{proof}
Since $M$ is free over the discrete valuation ring $O$, the image $L=EM$ of the idempotent $E$ is an $OC_p$-lattice,
and it is a direct summand of the permutation $OC_p$-lattice $M$.
By the characterization of $p$-permutation lattices recalled in Section \ref{sec:padic},
$L$ is therefore a $p$-permutation $OC_p$-lattice, and Lemma \ref{lem:trace-brauer} gives
\[
 \tr\bigl(x\mid K\otimes_O L\bigr)=\dim_\kappa L(C_p).
\]

It remains to compute this dimension.
The reduction $\overline E$ maps $C_p$-fixed vectors to $C_p$-fixed vectors and relative traces to relative traces,
because $E$ commutes with the action of $C_p$; therefore $E$ induces an idempotent endomorphism $E(C_p)$ of $M(C_p)$.
Since
\[
 M=L\oplus (1-E)M
\]
as $OC_p$-lattices, and since the Brauer quotient is additive on direct sums, this decomposition induces
\[
 M(C_p)=L(C_p)\oplus ((1-E)M)(C_p),
\]
and $E(C_p)$ is the projection onto the first summand.  Hence
\begin{equation}\label{eq:imageEC}
 \operatorname{im}E(C_p)=L(C_p).
\end{equation}

Under the identification \eqref{eq:MC-basis}, the matrix of $E(C_p)$ is $E[\mathcal F]\bmod\mathfrak m$.
To see this, apply $\overline E$ to $\overline e_w$ for a fixed point $w$.
Since $E$ commutes with the action of $x$ and $x(w)=w$,
the coefficients of $\overline E\,\overline e_w$ are constant on each $C_p$-orbit;
the component supported on an orbit of size $p$ is therefore a scalar multiple of the orbit sum,
i.e.\ a relative trace, and disappears in the Brauer quotient.
Thus only the coordinates at fixed points remain.
Combining this with \eqref{eq:imageEC},
\[
 \dim_\kappa L(C_p)=\rank_\kappa\bigl(E[\mathcal F]\bmod\mathfrak m\bigr).\qedhere
\]
\end{proof}

\end{document}
